\documentclass[10pt,a4paper]{amsart}
\usepackage[margin=19mm]{geometry}
\usepackage{amsmath,amssymb,mathtools}
\usepackage[scr=rsfs,cal=euler]{mathalfa}
\usepackage{xfrac}
\usepackage[unicode,hidelinks,bookmarks=false]{hyperref}
\newcommand{\ie}{i.e.\ }
\newcommand{\cf}{cf.\ }
\newcommand{\resp}{respectively\ }
\newcommand{\Subsection}{\subsection}
\providecommand{\nobreakdash}{\nobreak\hbox{-}}
\setlength{\emergencystretch}{2em}
\multlinegap0pt
\usepackage{booktabs}
\usepackage{caption}
\usepackage{algorithm}\usepackage{algpseudocode}
\usepackage[shortlabels]{enumitem}
\usepackage{csquotes}
\usepackage{amssymb}
\usepackage{float}
\newtheorem{theorem}{Theorem}
\title{Block Coordinate Descent Network Simplex Methods for Optimal Transport}
\author{Lingrui Li, Nobuo Yamashita}
\date{Source: arXiv:2506.21231v3 (CC BY 4.0)}
\begin{document}
\maketitle

\noindent\emph{Source and licence.} Block Coordinate Descent Network Simplex Methods for Optimal Transport. Version arXiv:2506.21231v3 (CC BY 4.0), distributed by arXiv under the Creative Commons Attribution 4.0 International licence, http://creativecommons.org/licenses/by/4.0/, as recorded in the arXiv licence metadata for this exact version and shown on the abstract page https://arxiv.org/abs/2506.21231v3. Full licence notice: \texttt{licenses/arxiv-2506.21231-CC-BY-4.0.txt}.\par\smallskip

\noindent\textit{Editorial scope.} Counted unit: the article's theorem thm:NS-global (source0.tex lines 478--497). The article's complete original proof of it is delivered with it (source0.tex lines 499--536, one \texttt{proof} environment of 38 lines). No mathematics is written, repaired or re-derived: the statement and the proof are the article's own lines, in the article's own notation, and no retained line is substituted or re-flowed.  No further source-local context is needed: every cross-reference of every retained line resolves inside the two delivered spans.  Nothing was omitted from any retained span, so the package carries no omission record.  No retained line contains a citation, so the wrapper carries no bibliography and no bibliography entry.  No retained line contains a figure environment, an image-inclusion command or drawing code, so no image is delivered and the package declares no asset dependency.  Editorial formatting only, in addition: an AMS \texttt{amsart} preamble that restates the article's own declarations and macros used by the retained lines, namely the environments \texttt{theorem} on separate counters, together with the standard packages those lines need.  The retained environments print in this document as Theorem 1 in order of appearance, whereas the article prints the same items with the numbering of its own full document.  The complete native source of the article as distributed in the arXiv e-print for this exact version is bundled unchanged as \texttt{sources/180540/source0.tex}.
\begin{theorem}\label{thm:NS-global}
Apply the NS method to the subproblem restricted to \(\mathcal H_{k}\), and set
\[
d^{k}\in\mathcal D\bigl(\mathbf{x}^{k};\mathcal H_{k}\bigr),
\qquad
\mathbf{x}^{k+1}=\mathbf{x}^{k}+d^{k}.
\]
Then:
\begin{enumerate}
  \item[\textnormal{(i)}]
  If \(\mathcal B_k\subseteq \mathcal H_k\) and \(\mathcal B^\ast\subseteq \mathcal H_k\),
  then \(\mathbf{x}^{k+1}\) is a globally optimal solution of the full OT problem.

  \item[\textnormal{(ii)}]
  Assume that the global optimal solution \(\mathbf{x}^\ast\) is \emph{unique}.
  If \(\mathcal{B}_{k}\not\subseteq\mathcal H_{k}\),
  then even when \(\mathcal B^\ast\subseteq \mathcal H_k\),
  the iterate \(\mathbf{x}^{k+1}\) cannot be globally optimal.
\end{enumerate}
\end{theorem}

\begin{proof}
\medskip
\noindent\textbf{(i)}
Because \(\mathcal B_k\subseteq \mathcal H_k\), the current basic feasible solution
\(\mathbf{x}^k\) is feasible for the restricted subproblem, and thus the NS method can be
initialized from \(\mathbf{x}^k\) without losing feasibility.
Let \(\mathbf{x}^{k+1}\) be an optimal solution of the restricted subproblem.
Then \(\mathbf{x}^{k+1}\) is feasible for the full OT problem (it satisfies \(A\mathbf{x}^{k+1}=b\) and
\(\mathbf{x}^{k+1}\ge 0\)), hence
\[
c^\top \mathbf{x}^{k+1}\ \ge\ c^\top \mathbf{x}^\ast.
\]
On the other hand, since \(\mathcal B^\ast\subseteq \mathcal H_k\), the basic feasible solution
\(\mathbf{x}^\ast\) (with nonbasic variables equal to zero) is feasible for the restricted subproblem as well.
By optimality of \(\mathbf{x}^{k+1}\) for the restricted subproblem, we have
\[
c^\top \mathbf{x}^{k+1}\ \le\ c^\top \mathbf{x}^\ast.
\]
Combining the two inequalities yields \(c^\top \mathbf{x}^{k+1}=c^\top \mathbf{x}^\ast\), so
\(\mathbf{x}^{k+1}\) attains the global optimal value of the full OT problem. Therefore
\(\mathbf{x}^{k+1}\) is globally optimal.

\medskip
\noindent\textbf{(ii)}
Assume \(\mathbf{x}^\ast\) is the unique globally optimal solution.
Let \(j\notin \mathcal H_k\) with \(x^k_j>0\). By the definition of the restricted subproblem, the \(j\)-th component cannot change:
\[
x^{k+1}_j \;=\; x^k_j + d^k_j \;=\; x^k_j \;>\; 0.
\]
Now suppose (for contradiction) that \(\mathbf{x}^{k+1}\) is globally optimal.
By uniqueness of the global optimizer, this would imply \(\mathbf{x}^{k+1}=\mathbf{x}^\ast\),
and in particular \(x^\ast_j=x^{k+1}_j>0\).

However, under the assumption \(\mathcal B^\ast\subseteq \mathcal H_k\), we have \(j\notin \mathcal B^\ast\),
and since \(\mathcal B^\ast\) is a basis of a basic feasible optimal solution, all variables outside
\(\mathcal B^\ast\) are nonbasic and therefore equal to zero. Thus \(x^\ast_j=0\), contradicting
\(x^\ast_j>0\). Consequently, \(\mathbf{x}^{k+1}\neq \mathbf{x}^\ast\), and \(\mathbf{x}^{k+1}\) cannot be globally optimal.
\end{proof}

\end{document}
